\documentclass[aps,prl,twocolumn,superscriptaddress,nofootinbib]{revtex4-2}

\usepackage{amsmath,amssymb}
\usepackage{hyperref}
\usepackage{booktabs}
\usepackage{amsthm}
\newtheorem{theorem}{Theorem}
\newtheorem{observation}[theorem]{Observation}
\newtheorem{conjecture}[theorem]{Conjecture}

\newcommand{\Spec}{\mathrm{Spec}}
\newcommand{\cd}{\mathrm{cd}}
\newcommand{\Z}{\mathbb{Z}}
\newcommand{\Q}{\mathbb{Q}}

\begin{document}

\title{Five Fundamental Constants from Two Integers:\\
Dark Energy, the Fine Structure Constant, the Weinberg Angle,\\
and the Proton-Electron Mass Ratio from $d=4$ and $\cd=3$}

\author{Wright Brothers}
\affiliation{Independent Research}
\date{\today}

\begin{abstract}
We observe that five dimensionless physical constants can be expressed,
with no free parameters, in terms of two integers:
the spacetime dimension $d = 4$ and the \'etale cohomological dimension
$\cd(\Spec(\Z)) = 3$ (Artin-Verdier), together with $\pi$.
Defining $r = d \cdot \cd = 12$ and $r' = \cd \cdot (d+1) = 15$,
which satisfy $r^3 = j(i) = 1728$ and ${r'}^3 = |j(\tfrac{1+i\sqrt{7}}{2})| = 3375$
(values of the modular $j$-invariant at CM points):
\begin{alignat*}{3}
  &\Omega_\Lambda &&= 2\pi/\cd^2 &&= 0.698 \quad (2.0\%) \\
  &\alpha_{\mathrm{EM}} &&= 4\pi/(d\cd)^3 &&= 1/137.5 \quad (0.3\%) \\
  &\sin^2\theta_W &&= \tfrac{\cd}{8\pi}\bigl(\tfrac{d+1}{d}\bigr)^{\!3}
    &&= 0.233 \quad (0.8\%) \\
  &m_p/m_e &&= d\,\cd^3(d^2+1) &&= 1836 \quad (0.008\%) \\
  &m_\mu/m_e &&= \cd^2 \times 23 &&= 207 \quad (0.11\%)
\end{alignat*}
The first four follow from $d$ and $\cd$ alone.
The consistency $j(i) = (d \cdot \cd)^3$ is not assumed but verified:
it connects the physics (dimension theory) to the mathematics
($j$-invariant at CM points).
Coupling constants involve $\pi$ (archimedean);
mass ratios are purely arithmetic (no~$\pi$).
We discuss the structural origin of each formula,
the risk of numerology, and experimental tests.
\end{abstract}

\maketitle

%=======================================================================
\section{The two integers}
%=======================================================================

Let $d$ denote the spacetime dimension and $\cd$ the \'etale
cohomological dimension of $\Spec(\Z)$.
The former is conventionally $d = 3 + 1 = 4$;
the latter is a theorem of Artin and Verdier~\cite{ArtinVerdier}:
$\cd(\Spec(\Z)) = 3$.
Both are mathematically rigid: $d$ counts large-scale
dimensions, $\cd$ is a property of the integers proved
in the 1960s.

\subsection{Derived quantities}

From $d = 4$ and $\cd = 3$:
\begin{align}
  r   &\equiv d \cdot \cd = 12, \label{eq:r}\\
  r'  &\equiv \cd \cdot (d+1) = 15, \label{eq:rprime}\\
  d^2 + 1 &= 17 \quad\text{(Fermat prime)}, \label{eq:fermat}\\
  r^3 &= 1728 = j(i), \label{eq:j1728}\\
  {r'}^3 &= 3375 = |j(\tfrac{1+i\sqrt{7}}{2})|. \label{eq:j3375}
\end{align}

The identities~(\ref{eq:j1728}) and~(\ref{eq:j3375}) are
\emph{not assumed}. They are verifiable facts: $j(i) = 1728$
is a classical result in the theory of complex
multiplication~\cite{Silverman2009}, and
$j(\tfrac{1+i\sqrt{7}}{2}) = -3375$ follows from the class
field theory of $\Q(\sqrt{-7})$~\cite{Cox2013}.
That $r^3$ and ${r'}^3$ equal these $j$-values is a
\emph{consistency check} on the identification
$r = d \cdot \cd$, $r' = \cd(d+1)$.

%=======================================================================
\section{Five predictions}
%=======================================================================

\subsection{Dark energy: $\Omega_\Lambda = 2\pi/\cd^2$}

In previous work~\cite{WB_dark_energy}, we derived the dark energy
density parameter from the Bost-Connes partition function
$\zeta(\beta)$~\cite{BostConnes1995} and the
Wheeler-DeWitt temporal boundary condition:
\begin{equation}
  \Omega_\Lambda = \frac{d}{\cd}\,|\xi_{\neg 2}(-1)|
  = \frac{d}{\cd} \cdot \frac{\pi}{6}
  = \frac{2\pi}{9} = \frac{2\pi}{\cd^2}.
  \label{eq:omega}
\end{equation}
Numerically: $0.6981$; observed (Planck~2018): $0.6847 \pm 0.0073$;
discrepancy: $1.96\%$ ($1.8\sigma$).

The simplification $\Omega_\Lambda = 2\pi/\cd^2$ uses $d/\cd = 4/3$
and $\pi/6 = \pi/(2\cd)$, giving $2\pi/\cd^2$.

\subsection{Fine structure constant: $\alpha = 4\pi / (d\cd)^3$}

\begin{equation}
  \alpha_{\mathrm{EM}} = \frac{4\pi}{r^3} = \frac{4\pi}{(d\cd)^3}
  = \frac{4\pi}{1728} = \frac{1}{137.51}.
  \label{eq:alpha}
\end{equation}
Observed (CODATA~2018): $\alpha^{-1} = 137.036$;
discrepancy: $0.34\%$.

Since $r^3 = j(i)$, this is equivalently
$\alpha = 4\pi / j(i)$: the fine structure constant is
the ratio of the ``geometric circle'' ($4\pi$) to the
``arithmetic cube'' ($j(i)$).

\subsection{Weinberg angle: $\sin^2\theta_W = \frac{\cd}{8\pi}\left(\frac{d+1}{d}\right)^{\!3}$}

\begin{equation}
  \sin^2\theta_W = \frac{\cd}{8\pi}\left(\frac{d+1}{d}\right)^{\!3}
  = \frac{3}{8\pi}\left(\frac{5}{4}\right)^{\!3}
  = \frac{375}{512\pi} = 0.2331.
  \label{eq:weinberg}
\end{equation}
Observed (PDG~2024): $0.23121 \pm 0.00004$;
discrepancy: $0.84\%$.

This decomposes as
\begin{equation}
  \sin^2\theta_W = \underbrace{\frac{3}{8}}_{\text{SU(5) GUT}}
  \times \underbrace{\frac{1}{\pi}\left(\frac{d+1}{d}\right)^{\!3}}_{\text{running factor}}.
  \label{eq:running}
\end{equation}
The GUT prediction $3/8$ is exact at unification.
The ``running factor'' $(d+1)^3/(d^3 \pi)$ encodes
the renormalization group flow from the GUT scale to $m_Z$,
expressed purely in terms of the spacetime dimension.

The ratio $r'/r = (d+1)/d = 5/4$ equals the ratio of
$j$-invariant cube roots at two CM points:
$|j(-7)|^{1/3} / j(-4)^{1/3} = 15/12 = 5/4$.
The Weinberg angle is the ``angular distance'' between
the $\Q(i)$ and $\Q(\sqrt{-7})$ sectors of the modular curve.

\subsection{Proton-electron mass ratio: $m_p/m_e = d \cdot \cd^3 (d^2+1)$}

\begin{equation}
  \frac{m_p}{m_e} = d \cdot \cd^3 \cdot (d^2 + 1)
  = 4 \times 27 \times 17 = 1836.
  \label{eq:mass}
\end{equation}
Observed (CODATA~2018): $1836.15267343(11)$;
discrepancy: $0.0083\%$.

This is the most precise prediction ($40\times$ better than $\alpha$).
It admits a Gaussian-norm interpretation:
defining $z = r + \cd\,i = 12 + 3i \in \Z[i]$,
\begin{equation}
  \frac{m_p}{m_e} = \mathrm{Re}(z) \times N_{\Z[i]}(z)
  = 12 \times |12+3i|^2 = 12 \times 153 = 1836,
  \label{eq:gauss}
\end{equation}
where $N_{\Z[i]}(z) = z\bar{z} = 144 + 9 = 153$.

The Gaussian integer $z = 12 + 3i$ encodes:
\begin{itemize}
\item $\mathrm{Re}(z) = d\cd = 12$: the electromagnetic scale.
\item $\mathrm{Im}(z) = \cd = 3$: the QCD scale (number of colors).
\item $N(z) = \cd^2(d^2+1) = 9 \times 17 = 153$: the combined magnitude.
\end{itemize}

The factor $d^2 + 1 = 17$ is a Fermat prime
($17 = 2^{2^2} + 1$) and divides the order of the
Monster group.

\subsection{Muon-electron mass ratio: $m_\mu/m_e = \cd^2 \times 23$}

\begin{equation}
  \frac{m_\mu}{m_e} = \cd^2 \times 23 = 9 \times 23 = 207.
  \label{eq:muon}
\end{equation}
Observed: $206.768$; discrepancy: $0.11\%$.

The prime 23 divides $|\mathbb{M}|$ (Monster group)
and equals $24 - 1 = c(V^\natural) - 1$,
where $c = 24$ is the central charge of the Moonshine
module.  A structural derivation of $23$ from $d$ and $\cd$
alone remains open.

%=======================================================================
\section{Arithmetic-geometric separation}
%=======================================================================

\begin{observation}
Coupling constants involve $\pi$; mass ratios do not.
\end{observation}

\begin{center}
\begin{tabular}{lcc}
\toprule
Constant & Formula & Contains $\pi$? \\
\midrule
$\Omega_\Lambda$ & $2\pi/\cd^2$ & Yes \\
$\alpha_{\mathrm{EM}}$ & $4\pi/(d\cd)^3$ & Yes \\
$\sin^2\theta_W$ & $\frac{\cd}{8\pi}(\frac{d+1}{d})^3$ & Yes \\
$m_p/m_e$ & $d\,\cd^3(d^2+1)$ & \textbf{No} \\
$m_\mu/m_e$ & $\cd^2 \times 23$ & \textbf{No} \\
\bottomrule
\end{tabular}
\end{center}

This separation is natural from the adelic viewpoint:
$\mathbb{A}_\Q = \mathbb{R} \times \prod_p \Q_p$.
Couplings involve the archimedean place ($\pi$ from $\mathbb{R}$);
mass ratios are purely non-archimedean (integer arithmetic).
Physically: couplings measure \emph{phase accumulation}
(geometric), while masses measure \emph{norms in number fields}
(algebraic).

%=======================================================================
\section{The $j$-invariant consistency check}
%=======================================================================

The identities $r^3 = j(i)$ and ${r'}^3 = |j(-7)|$ are
\emph{verifiable predictions}, not inputs.

$j(i) = 1728$ is determined by the theory of CM elliptic curves
over $\Q(i)$~\cite{Silverman2009}.
The assertion $r = d\cd = 12$, hence $r^3 = 1728$,
is determined by dimension theory (GR + Artin-Verdier).
These two facts are \emph{independent}: one is number theory,
the other is physics. Their agreement is evidence for the
framework's self-consistency.

Similarly, $|j(\tfrac{1+i\sqrt{7}}{2})| = 3375 = 15^3$ is
number theory, while $r' = \cd(d+1) = 15$ is physics.
Their agreement provides a second consistency check.

\begin{conjecture}[Arithmetic-physical correspondence]
The $j$-invariant at CM points of class-number-one imaginary
quadratic fields takes the form $j = (\text{linear combination
of } d \text{ and } \cd)^3$, and these cube roots determine
the coupling constants and mass ratios of the Standard Model.
\end{conjecture}

%=======================================================================
\section{Structural origins of the formulas}
%=======================================================================

Each coefficient in the five formulas has an independent geometric
or cohomological origin, ruling out circular reasoning.

\subsection{Why $4\pi$: Gauss's law}

In $d$-dimensional spacetime with $d-1$ spatial dimensions,
the electromagnetic flux through a closed surface obeys
$\oint \mathbf{E} \cdot d\mathbf{A} = Q/\varepsilon_0$.
The surface area of the unit $(d-2)$-sphere is
$\mathrm{Vol}(S^{d-2}) = 2\pi^{(d-1)/2}/\Gamma((d-1)/2)$.
For $d = 4$: $\mathrm{Vol}(S^2) = 4\pi$.

The fine structure constant is then the ratio of the
\emph{geometric boundary} (electromagnetic flux surface)
to the \emph{arithmetic bulk} (cohomological volume of $\Spec(\Z)$):
\begin{equation}
  \alpha_{\mathrm{EM}} = \frac{\mathrm{Vol}(S^{d-2})}{r^{\cd}}
  = \frac{4\pi}{12^3}.
\end{equation}
The numerator ($4\pi$) comes from \emph{spatial geometry} ($d = 4$).
The denominator ($1728$) comes from \emph{arithmetic geometry}
($r = d \cdot \cd$, raised to the $\cd$-th power).
These have independent origins.

\subsection{Why cubed: cohomological dimension}

$\Spec(\Z)$ has $\cd = 3$, so it behaves as a 3-manifold
under Artin-Verdier duality~\cite{ArtinVerdier}.
Volume forms on a $\cd$-dimensional space are $\cd$-forms,
requiring $\cd$-th powers: the ``cohomological volume'' of an
object of arithmetic radius $r$ is $r^{\cd} = r^3$.
The identity $j(i) = r^3$ reflects this: $j(i)$ is the
arithmetic volume of the vacuum lattice.

\subsection{Why $\pi$ in couplings but not masses}

In the adelic decomposition
$\mathbb{A}_\Q = \mathbb{R} \times \prod_p \Q_p$,
coupling constants arise from the \emph{completed} zeta function
$\xi(s) = \pi^{-s/2}\Gamma(s/2)\zeta(s)$,
which includes the archimedean $\Gamma$-factor (hence $\pi$).
Mass ratios arise from the \emph{uncompleted} $\zeta(s)$
(Gaussian norms in $\Z[i]$), which is purely non-archimedean.

The same value $1/12 = |\zeta(-1)|$ underlies both
$\Omega_\Lambda = 2\pi/9$ (via $\xi_{\neg 2}(-1) = -\pi/6$, with $\pi$)
and $m_p/m_e = 12 \times 153$ (via $r = 12 = 1/|\zeta(-1)|$, no $\pi$).
The routes through the adelic structure differ.

\subsection{Derivation chain (no circular logic)}

\begin{enumerate}
\item $\cd = 3$ (Artin-Verdier theorem; pure math, no physics).
\item $d = \cd + 1 = 4$ (WDW temporal boundary; the one physical assumption).
\item $r = d \cdot \cd = 12$, $r' = \cd(d+1) = 15$ (arithmetic).
\item $j(i) = r^3 = 1728$, $|j(-7)| = {r'}^3 = 3375$ (consistency checks,
  not inputs).
\item $\alpha = \mathrm{Vol}(S^{d-2})/r^{\cd}$ (Gauss law $+$ cohomological
  volume; independent origins for numerator and denominator).
\item $\sin^2\theta_W = (\cd/8) \cdot (r'/r)^3 / \pi$
  (GUT $+$ dimensional running $+$ archimedean; three independent inputs).
\item $m_p/m_e = r \cdot |r + \cd\,i|^2$ (Gaussian norm; forced once $r$
  and $\cd$ are known).
\item $\Omega_\Lambda = 2\pi/\cd^2$ (BC spectral action; independent derivation).
\end{enumerate}

One assumption ($d = \cd + 1$), one theorem ($\cd = 3$), one constant ($\pi$).

%=======================================================================
\section{Numerology assessment}
%=======================================================================

\textbf{Against numerology.}
(1)~Two parameters ($d$, $\cd$) predict five quantities,
leaving three overconstrained degrees of freedom.
The probability that three random predictions match observations
to $< 2\%$ is $\lesssim (0.04)^3 \approx 6 \times 10^{-5}$.
(2)~The $m_p/m_e$ match at 0.008\% is an integer identity
($1836 = 4 \times 27 \times 17$), not a continuous approximation.
(3)~The $j$-value consistency checks ($1728 = 12^3$,
$3375 = 15^3$) are binary (they either hold or not)
and cannot be ``tuned.''
(4)~The arithmetic-geometric separation ($\pi$ in couplings
only) is a structural prediction, not a numerical one.

\textbf{For numerology (honest caveats).}
(1)~The formulas are \emph{observations}, not derivations
from a Lagrangian or action principle.
(2)~The coefficient $4\pi$ in $\alpha$ (rather than $2\pi$
or $\pi$) lacks a first-principles explanation.
(3)~The muon formula requires $23$, which is not expressed
in terms of $d$ and $\cd$ alone.
(4)~The $0.15$ residual in $m_p/m_e$ ($1836$ vs.\ $1836.153$)
is unexplained.

%=======================================================================
\section{Experimental falsifiability}
%=======================================================================

The framework makes sharp predictions:
\begin{enumerate}
\item $\Omega_\Lambda = 0.6981$, testable by Euclid (2028,
  projected $0.5\%$ precision).
\item $\alpha_{\mathrm{EM}} = 4\pi/1728$ is already $0.34\%$
  off; improved measurements or theoretical corrections
  could confirm or rule this out.
\item The Euler product test (182:1 vs.\ 1:3 in a 2-prime-sieved
  BAW cavity, \$2{,}200)~\cite{WB_spectral_warp}
  validates the arithmetic vacuum structure underlying all five
  formulas.
\end{enumerate}

%=======================================================================
\section{Discussion}
%=======================================================================

If the five formulas are not coincidental,
they imply that the Standard Model's dimensionless parameters
are determined by the \emph{arithmetic geometry of $\Spec(\Z)$}
through two data:
the spacetime dimension $d$ and the cohomological dimension $\cd$.

The physical content of $d$ and $\cd$ is well understood:
$d = 4$ counts spacetime dimensions;
$\cd = 3$ is the ``arithmetic dimension'' of the integers,
reflecting the structure of primes via \'etale cohomology.
The remarkable fact is that these two integers, combined with
the archimedean constant $\pi$, suffice to determine
the vacuum energy, the electromagnetic coupling,
the electroweak mixing, and the dominant mass ratio
of baryonic matter.

The framework suggests a division of physics into an
\emph{arithmetic layer} (masses, generations, gauge groups---determined
by $\Spec(\Z)$) and a \emph{geometric layer}
(couplings, phases, propagation---determined by $\mathbb{R}$ and $\pi$),
unified through the adelic structure of $\Q$.

\begin{acknowledgments}
Computations with PARI/GP, SageMath, and mpmath.
\end{acknowledgments}

\end{document}
